\chapter{Segmentaci\'on}
\label{cap:segmentacion}

La segmentación divide una tabla en fragmentos lógicos mediante vistas. A diferencia del particionamiento, la tabla no se modifica físicamente: sus datos permanecen en un solo lugar y cada fragmento es una vista que selecciona una parte de sus filas o de sus columnas. Los usuarios consultan el fragmento que necesitan como si fuera una tabla, y el gestor obtiene sus datos de la tabla original en cada consulta.

Como se indicó en el capítulo~\ref{cap:multiinstancias}, la segmentación se implementa en la segunda instancia de cada gestor, junto con el particionamiento.

\section{Dise\~no de la segmentaci\'on}

\paragraph{Tipos de fragmentación.} El esquema aplica los dos tipos de fragmentación, cada uno sobre la tabla en la que tiene sentido (tabla~\ref{tab:fragmentacion}).

\begin{table}[H]
	\centering
	\small
	\caption{Tipos de fragmentación aplicados}
	\label{tab:fragmentacion}
	\begin{tabular}{@{}L{2.4cm} L{2.2cm} L{5.2cm} L{6cm}@{}}
		\toprule
		\textbf{Tipo} & \textbf{Tabla} & \textbf{Criterio} & \textbf{Cada fragmento contiene} \\
		\midrule
		Horizontal & \nem{inversion} & Estado de la inversión & Todas las columnas de las inversiones que se encuentran en un estado \\
		Vertical & \nem{bitacora} & Uso de las columnas & Todos los registros, con un subconjunto de las columnas y la llave primaria \\
		\bottomrule
	\end{tabular}
\end{table}

\paragraph{Reglas de corrección.} Una fragmentación es correcta si cumple tres reglas, que el diseño de los fragmentos de este esquema garantiza:

\begin{itemize}[leftmargin=*]
	\item \textbf{Completitud:} todo dato de la tabla aparece en algún fragmento. En la fragmentación horizontal, toda fila pertenece a un fragmento; en la vertical, toda columna.
	\item \textbf{Disyunción:} ningún dato aparece en más de un fragmento. En la fragmentación vertical, la excepción es la llave primaria, que se repite en todos los fragmentos para poder reunirlos.
	\item \textbf{Reconstrucción:} la tabla original puede obtenerse a partir de sus fragmentos: con la unión de los fragmentos horizontales (\nem{UNION ALL}) o con la reunión de los fragmentos verticales por su llave primaria (\nem{JOIN}).
\end{itemize}

\subsection{Fragmentación horizontal de \nem{inversion}}

La tabla \nem{inversion} se divide en tres fragmentos, uno por cada estado de una inversión definido en las reglas de negocio (tabla~\ref{tab:fragmentos-inversion}).

\begin{table}[H]
	\centering
	\small
	\caption{Fragmentos horizontales de \nem{inversion}}
	\label{tab:fragmentos-inversion}
	\begin{tabular}{@{}L{3.6cm} L{7.4cm} L{4.8cm}@{}}
		\toprule
		\textbf{Vista} & \textbf{Condición} & \textbf{Uso principal} \\
		\midrule
		\nem{inversion\us vigente} & Sin liquidación y con vencimiento igual o posterior a la fecha actual & Seguimiento del ejecutivo y del gerente \\
		\nem{inversion\us vencida} & Sin liquidación y con vencimiento anterior a la fecha actual & Cobro de adeudos de la tesorería \\
		\nem{inversion\us liquidada} & Con liquidación registrada & Revisión del auditor \\
		\bottomrule
	\end{tabular}
\end{table}

Las condiciones son excluyentes y cubren todos los casos: una inversión tiene o no tiene liquidación, y si no la tiene, su vencimiento es anterior a la fecha actual o no lo es. Por ello, la fragmentación es completa y disjunta por construcción.

\paragraph{Por qué vistas y no particiones.} El estado de una inversión no se almacena en ninguna columna: depende de la fecha actual y de la existencia de un registro en otra tabla, \nem{liquidacion}. Una inversión vigente pasa a estar vencida el día siguiente a su vencimiento sin que nadie modifique un solo registro. Por ello, el estado no puede servir como clave de partición: el gestor ubica cada fila en una partición según el valor almacenado en una columna de la propia tabla, y lo hace al insertarla o modificarla. Agregar una columna de estado obligaría a actualizarla diariamente con un proceso externo y repetiría información que ya se deduce de los datos. Una vista, en cambio, evalúa la condición en cada consulta, por lo que su contenido siempre corresponde al estado actual.

\paragraph{Contenido variable.} Como consecuencia, el contenido de las vistas cambia con el paso del tiempo. Con los datos de prueba, las 40 inversiones se reparten en 15 vigentes, 10 vencidas y 15 liquidadas. Mientras no se registren nuevas liquidaciones, estos conteos se mantienen hasta el 19 de enero de 2027, fecha de vencimiento de la inversión 30; al día siguiente, esa inversión pasa al fragmento de vencidas. Las reglas de corrección, en cambio, se cumplen en cualquier fecha.

\subsection{Fragmentación vertical de \nem{bitacora}}

Las columnas de \nem{bitacora} se usan de dos formas distintas. Las que identifican el evento (quién, cuándo, desde dónde, en qué tabla y con qué operación) se consultan constantemente para revisar la actividad del sistema. Las que contienen la sentencia y los valores anteriores y nuevos solo se consultan al revisar un cambio en particular, y son además las más pesadas, de tipo \nem{TEXT} y \nem{JSON}. La tabla se divide en dos fragmentos según ese uso (tabla~\ref{tab:fragmentos-bitacora}).

\begin{table}[H]
	\centering
	\small
	\caption{Fragmentos verticales de \nem{bitacora}}
	\label{tab:fragmentos-bitacora}
	\begin{tabular}{@{}L{3.4cm} L{8.2cm} L{4.2cm}@{}}
		\toprule
		\textbf{Vista} & \textbf{Columnas} & \textbf{Uso principal} \\
		\midrule
		\nem{bitacora\us resumen} & \nem{bitacora\us num}, \nem{bitacora\us fecha}, \nem{bitacora\us usuario}, \nem{bitacora\us tabla}, \nem{bitacora\us operacion}, \nem{bitacora\us ip}, \nem{bitacora\us sesion} & Revisión de la actividad: quién hizo qué y cuándo \\
		\nem{bitacora\us detalle} & \nem{bitacora\us num}, \nem{bitacora\us fecha}, \nem{bitacora\us query}, \nem{bitacora\us dato\us anterior}, \nem{bitacora\us dato\us nuevo} & Revisión del contenido de un cambio \\
		\bottomrule
	\end{tabular}
\end{table}

Ambos fragmentos conservan la llave primaria, (\nem{bitacora\us num}, \nem{bitacora\us fecha}), que permite reunirlos para reconstruir cada registro. Cada una de las demás columnas aparece en un solo fragmento, por lo que la fragmentación vertical también es completa y disjunta por construcción. El auditor localiza primero los eventos en \nem{bitacora\us resumen} y después consulta el detalle de los que le interesan.

\paragraph{Por qué vistas y no tablas separadas.} Dividir la bitácora en dos tablas físicas obligaría a cada trigger a escribir en ambas, y en MySQL no podría haber una llave foránea entre ellas, porque las tablas particionadas no admiten llaves foráneas. Con vistas, los triggers siguen escribiendo un solo registro en una sola tabla, y la fragmentación no modifica la forma de auditar descrita en el capítulo~\ref{cap:bitacora}. Las vistas tampoco dependen de que la tabla esté particionada: el gestor resuelve las particiones al consultarlas.

\subsection{Uso y privilegios de las vistas}

Las vistas son de solo lectura para los usuarios del sistema: las altas y los cambios se siguen haciendo sobre las tablas, con los privilegios del capítulo~\ref{cap:implementacion}. Cada vista se ejecuta con los privilegios de su dueño, el DBA, por lo que otorgar el privilegio \nem{SELECT} sobre una vista da acceso a sus datos aunque el usuario no tenga acceso a la tabla. Por ello, los privilegios sobre las vistas siguen la misma matriz que los de las tablas:

\begin{itemize}[leftmargin=*]
	\item Los fragmentos de \nem{inversion} se otorgan al gerente, al ejecutivo y a la tesorería, que ya pueden consultar la tabla completa.
	\item Los fragmentos de \nem{bitacora} no se otorgan a ningún rol de operación, ya que solo el auditor puede consultar la bitácora.
	\item El auditor y el usuario de respaldo los reciben sin otorgamientos adicionales: en PostgreSQL, por los privilegios por defecto definidos en el capítulo~\ref{cap:implementacion}; en MySQL, porque sus privilegios abarcan toda la base de datos (\nem{inversiones.*}).
\end{itemize}

Las vistas se definen con \nem{SELECT i.*} en los fragmentos horizontales, ya que estos conservan todas las columnas. Ambos gestores expanden el asterisco al crear la vista, por lo que, si en el futuro se agrega una columna a \nem{inversion}, las vistas deben volver a crearse para incluirla.

\section{Diagrama relacional con segmentaci\'on}

Las figuras~\ref{fig:segmentacion-horizontal} y~\ref{fig:segmentacion-vertical} muestran las vistas que se derivan de cada tabla. Las tablas se dibujan con borde continuo y encabezado oscuro, y las vistas con borde punteado y encabezado claro. Las flechas continuas indican de qué tabla obtiene sus datos cada vista; las punteadas, la tabla que se consulta para evaluar su condición. Las demás tablas y relaciones del esquema no cambian respecto al diagrama relacional (figura~\ref{fig:diagrama-er}).

\captura{segmentacion_horizontal.jpg}{Fragmentación horizontal de \nem{inversion}}{segmentacion-horizontal}
Las flechas continuas indican el origen de los datos; las punteadas, la tabla que se consulta para evaluar la condición.

\captura[0.85\textwidth]{segmentacion_vertical.jpg}{Fragmentación vertical de \nem{bitacora}}{segmentacion-vertical}
Ambas vistas conservan la llave primaria, que permite reconstruir cada registro.

\section{PostgreSQL}

El DBA crea las vistas en la segunda instancia, con el rol \nem{rol\us admin\us bd} activo para que el rol sea su dueño, igual que el de las tablas. En la misma secuencia otorga los privilegios sobre los fragmentos de \nem{inversion}.

\begin{verbatim}
	-- Ejecutado por dlandero en la segunda instancia, base de datos inversiones
	SET ROLE rol_admin_bd;

	-- Fragmentación horizontal de inversion, por estado
	CREATE VIEW inversion_vigente AS
	SELECT i.*
	FROM inversion i
	WHERE i.inversion_fecha_vencimiento >= CURRENT_DATE
	  AND NOT EXISTS (SELECT 1 FROM liquidacion l
	                  WHERE l.inversion_num = i.inversion_num);

	CREATE VIEW inversion_vencida AS
	SELECT i.*
	FROM inversion i
	WHERE i.inversion_fecha_vencimiento < CURRENT_DATE
	  AND NOT EXISTS (SELECT 1 FROM liquidacion l
	                  WHERE l.inversion_num = i.inversion_num);

	CREATE VIEW inversion_liquidada AS
	SELECT i.*
	FROM inversion i
	WHERE EXISTS (SELECT 1 FROM liquidacion l
	              WHERE l.inversion_num = i.inversion_num);

	-- Fragmentación vertical de bitacora
	CREATE VIEW bitacora_resumen AS
	SELECT bitacora_num, bitacora_fecha, bitacora_usuario, bitacora_tabla,
	       bitacora_operacion, bitacora_ip, bitacora_sesion
	FROM bitacora;

	CREATE VIEW bitacora_detalle AS
	SELECT bitacora_num, bitacora_fecha, bitacora_query,
	       bitacora_dato_anterior, bitacora_dato_nuevo
	FROM bitacora;

	-- Privilegios
	GRANT SELECT ON inversion_vigente, inversion_vencida, inversion_liquidada
	  TO rol_gerente, rol_ejecutivo, rol_tesoreria;
\end{verbatim}

\section{MySQL}

Las vistas tienen la misma definición que en PostgreSQL. Se crean con \nem{DEFINER = CURRENT\us USER} y \nem{SQL SECURITY DEFINER}, de modo que se ejecutan con los privilegios de la cuenta del DBA, igual que los procedimientos y triggers del esquema. Como MySQL solo permite otorgar privilegios sobre un objeto por sentencia, cada fragmento de \nem{inversion} requiere su propio \nem{GRANT}.

\begin{verbatim}
	-- Ejecutado por dlandero en la segunda instancia, base de datos inversiones

	-- Fragmentación horizontal de inversion, por estado
	CREATE DEFINER = CURRENT_USER SQL SECURITY DEFINER VIEW inversion_vigente AS
	SELECT i.*
	FROM inversion i
	WHERE i.inversion_fecha_vencimiento >= CURRENT_DATE
	  AND NOT EXISTS (SELECT 1 FROM liquidacion l
	                  WHERE l.inversion_num = i.inversion_num);

	CREATE DEFINER = CURRENT_USER SQL SECURITY DEFINER VIEW inversion_vencida AS
	SELECT i.*
	FROM inversion i
	WHERE i.inversion_fecha_vencimiento < CURRENT_DATE
	  AND NOT EXISTS (SELECT 1 FROM liquidacion l
	                  WHERE l.inversion_num = i.inversion_num);

	CREATE DEFINER = CURRENT_USER SQL SECURITY DEFINER VIEW inversion_liquidada AS
	SELECT i.*
	FROM inversion i
	WHERE EXISTS (SELECT 1 FROM liquidacion l
	              WHERE l.inversion_num = i.inversion_num);

	-- Fragmentación vertical de bitacora
	CREATE DEFINER = CURRENT_USER SQL SECURITY DEFINER VIEW bitacora_resumen AS
	SELECT bitacora_num, bitacora_fecha, bitacora_usuario, bitacora_tabla,
	       bitacora_operacion, bitacora_ip, bitacora_sesion
	FROM bitacora;

	CREATE DEFINER = CURRENT_USER SQL SECURITY DEFINER VIEW bitacora_detalle AS
	SELECT bitacora_num, bitacora_fecha, bitacora_query,
	       bitacora_dato_anterior, bitacora_dato_nuevo
	FROM bitacora;

	-- Privilegios
	GRANT SELECT ON inversiones.inversion_vigente
	  TO rol_gerente, rol_ejecutivo, rol_tesoreria;
	GRANT SELECT ON inversiones.inversion_vencida
	  TO rol_gerente, rol_ejecutivo, rol_tesoreria;
	GRANT SELECT ON inversiones.inversion_liquidada
	  TO rol_gerente, rol_ejecutivo, rol_tesoreria;
\end{verbatim}

